
% =============================================================================
% Project : THE PHILOSOPHY OF VIABILITY
% Book    : Book I — Ontology of Viability
% Book ID : VS.B01
% Part    : Part I — Identity and Manifestation
% Part ID : VS.B01.P01
% Chapter : Chapter 08 — Comparative Philosophical Examination
% Chapter ID : VS.B01.P01.C08
%
% File :
%   4_ASSETS/tables/
%   VS.B01.P01.C08.S16.SS18.P004.TL01.ibn_arabi_claims_and_comparative_limits.tex
%
% Asset : TL01
% Structural Identifier :
%   VS.B01.P01.C08.S16.SS18.P004.TL01
%
% Purpose :
%   To distinguish the reconstructed claims of Ibn Arabi
%   from their logical forms and epistemic statuses,
%   while recording their comparative limits for Book I.
%
% Dependencies :
%   longtable, booktabs, array.
%
% =============================================================================

\begingroup
\small
\renewcommand{\arraystretch}{1.2}
\setlength{\tabcolsep}{4pt}
\setlength{\LTcapwidth}{\linewidth}

\def\colw#1{%
    \dimexpr(\linewidth-8\tabcolsep)*#1/100\relax
}

\def\rowsep{%
    \specialrule{0.2pt}{3pt}{3pt}
}

\def\tablehead{%
    \toprule
    \textbf{Concept or Claim} &
    \textbf{Akbarian Reconstruction} &
    \textbf{Logical Form} &
    \textbf{Epistemic Status} &
    \textbf{Comparative Limit for Book I}
    \\
    \midrule
}

\def\grouprow#1{%
    \multicolumn{5}{@{}l}{\footnotesize\textsc{#1}}\\
    \cmidrule{1-5}
    \addlinespace[1pt]
}

\begin{longtable}{@{}
    >{\raggedright\arraybackslash}p{\colw{14}}
    >{\raggedright\arraybackslash}p{\colw{29}}
    >{\raggedright\arraybackslash}p{\colw{14}}
    >{\raggedright\arraybackslash}p{\colw{17}}
    >{\raggedright\arraybackslash}p{\colw{26}}
@{}}

\caption{
    Ibn Arabi's Principal Ontological Claims and Their
    Comparative Limits
}
\label{tab:VS.B01.P01.C08.S16.SS18.P004.TL01}
\\

\tablehead
\endfirsthead

\multicolumn{5}{@{}l}{%
    \footnotesize\textit{
        Table~\thetable\ (continued):
        Ibn Arabi's Principal Ontological Claims
    }
}
\\

\tablehead
\endhead

\multicolumn{5}{@{}r}{%
    \footnotesize\textit{Continued on next page}
}
\\
\endfoot

\bottomrule
\endlastfoot

\grouprow{Unity and the Real}

\textbf{Tawhid}
&
The affirmation of the unity of the Real and the
dependence of possible existence upon it.
&
Thesis
&
Argued/Derived Claims
&
The unity of the Real is not equivalent to the
unity or singularity of Identity.
\\
\rowsep

\textbf{Wahdat al-wujud}
&
A later expression used in the history of interpretation;
not Ibn Arabi's own technical designation.
&
Historical interpretation
&
Argued/Derived Claims
&
The later label must not be projected without qualification
onto the primary texts.
\\
\rowsep

\textbf{Nondelimitation}
&
The unrestricted character of the Real in relation
to determinate and delimited existence.
&
Definition (historical reconstruction)
&
Argued/Derived Claims
&
Nondelimitation is not equivalent to Invariance,
Immutability, or Identity.
\\
\rowsep

\grouprow{Names and manifestation}

\textbf{Divine names}
&
Relational meanings and attributes through which
the Real is considered under differentiated respects.
&
Thesis
&
Argued/Derived Claims
&
Divine names are not identical with Determination
in Book I.
\\
\rowsep

\textbf{Tajalli}
&
The self-disclosure of the Real through determinate
forms and relations.
&
Thesis
&
Argued/Derived Claims
&
Tajalli is not equivalent to generic Manifestation.
\\
\rowsep

\textbf{Preparedness}
&
The specific capacity of a locus to receive and
display determinate manifestations.
&
Thesis
&
Argued/Derived Claims
&
Preparedness is not Identity or a complete
account of the conditions of manifestation in Book I.
\\
\rowsep

\grouprow{Determinate realities}

\textbf{A'yan thabita}
&
Determinate realities considered in relation to
divine knowledge without independent existence
in themselves.
&
Thesis
&
Argued/Derived Claims
&
Fixed entities must not be identified with Identity
or with independently existing Platonic archetypes.
\\
\rowsep

\textbf{A'yan mawjuda}
&
The same determinate realities considered through
their manifested or delimited existence.
&
Thesis
&
Argued/Derived Claims
&
Their relation to fixed entities does not establish
a Book I theory of identity through transformation.
\\
\rowsep

\grouprow{Intermediate realities}

\textbf{Barzakh}
&
A limit or isthmus that simultaneously separates
and joins two sides.
&
Definition (historical reconstruction)
&
Argued/Derived Claims
&
The Barzakh is not equivalent to any one of
Book I's three mediating hubs.
\\
\rowsep

\textbf{Imaginal realm}
&
An ontologically significant intermediate domain
between spiritual and corporeal forms.
&
Thesis
&
Conditionally Supported Claims
&
It is not automatically a separate realm
within Book I's architecture.
\\
\rowsep

\grouprow{Renewal and comprehensive reality}

\textbf{Non-repetition of tajalli}
&
Self-disclosure is not repeated in precisely
the same manner; creation is continuously renewed.
&
Thesis
&
Argued/Derived Claims
&
Non-repetition of manifestations does not
establish the loss or plurality of Identity.
\\
\rowsep

\textbf{Perfect Human}
&
A comprehensive human reality related to
the divine and cosmic dimensions.
&
Thesis
&
Conditionally Supported Claims
&
The concept has no direct equivalent in the
foundational ontology of Book I.
\\
\rowsep

\textbf{Five Divine Presences}
&
A systematic scheme developed in the later
Akbarian tradition, especially through Qunawi
and subsequent interpreters.
&
Historical interpretation
&
Conditionally Supported Claims
&
The later scheme must not be attributed
unqualifiedly to Ibn Arabi.
\\
\rowsep

\grouprow{Comparative status}

\textbf{Being and Identity}
&
Book I maintains Analytical Distinction,
Ontological Inseparability, and Absence
of Ontological Priority.
&
Foundational Commitment
&
Foundationally Adopted Commitments
&
These commitments are not derived from
Ibn Arabi's account.
\\
\rowsep

\textbf{Identity as fixed entity}
&
The identification of Book I Identity with
the Akbarian fixed entity is rejected.
&
Rejected Thesis
&
Rejected/Superseded Claims
&
The two concepts perform different
ontological and explanatory functions.
\\
\rowsep

\textbf{Comparative synthesis}
&
The two systems address overlapping questions
through distinct metaphysical architectures.
&
Thesis
&
Provisional Constructions
&
Conceptual resemblance does not establish
equivalence, derivation, or historical influence.
\\
\rowsep

\textbf{Plurality of identities}
&
The comparison does not settle whether
multiple identities obtain and under
which conditions.
&
Open Problem
&
Unresolved/Open Problems
&
No conclusion follows merely from the
plurality of names or manifestations.
\\

\end{longtable}

\endgroup

% =============================================================================
% END OF TABLE TL01
% =============================================================================
